\documentclass[10pt, titlepage]{article}

\usepackage[margin=1in]{geometry}
\usepackage{fancyhdr}
\usepackage{amsfonts}
\usepackage{amsmath}
\usepackage{amssymb}
\usepackage{booktabs}
%\usepackage{tikz}
\usepackage{float}
\usepackage{bm}
\usepackage{pgfplots}
\usepackage{array}
\usepackage{nicematrix}

\title{\normalsize \textbf{Double} PLUS Session Worksheet \\[1em] \Large \textit{MATH 1554:} Linear Algebra}
\author{\textit{Exam 3 Review} \\ by \textit{Ahmed Mohammed}, \textit{Nora Butler}, \textit{Sri Palaniappan}}
\date{Fall 2023}

% ------- Set up header -------
\renewcommand{\sectionmark}[1]{\markright{#1}} % from https://tex.stackexchange.com/a/317393/227665
\pagestyle{fancy}
\fancyhead[L]{\textbf{Double 3}, Fall 2023}
\fancyhead[C]{\textit{\nouppercase{\rightmark}}}
%\fancyhead[C]{\textit{Linear Transformations (1.8)}}
\fancyhead[R]{\textbf{\thepage}}
% -------------------------------------------------------

% ------- Set up question environment -------
\newcounter{question}
\setcounter{question}{1}
\newenvironment{newquestion}
{
	\noindent\textbf{Question \thequestion}
}
{
\stepcounter{question}
\par
}
% -------------------------------------------------------

% ------- Set up Main Idea box environment -------
\newcounter{mainideas}
\setcounter{mainideas}{1}
\newcommand{\mainidea}{\textbf{Main Idea \#\themainideas \stepcounter{mainideas}}}
\setlength\fboxsep{1em}
%\newcommand{\midea}[2]{\begin{center}\fbox{\parbox{#1}{\centering \mainidea #2}} \end{center}} % Could be environment but I couldn't bother

% The way this works is that the savebox collects whatever was between the minipage and the end of it, then in the last line of the definition of the environment
% we regurgigate whatever was collected into an fbox which we center.
\newsavebox{\mybox}
\newenvironment{mideea}[1]
{%
\begin{lrbox}{\mybox}%
	\begin{minipage}{#1\textwidth}
		{\centering \mainidea \par}

}
{\end{minipage}\end{lrbox}\begin{center}\fbox{\usebox{\mybox}}\end{center}}
% -------------------------------------------------------


% ------- Set up math stuff -------
\renewcommand{\vec}[1]{\bm{#1}}
\DeclareMathOperator{\vspan}{Span}
\DeclareMathOperator{\row}{row}
\DeclareMathOperator{\col}{col}
\DeclareMathOperator{\colspace}{Col}
\DeclareMathOperator{\rowspace}{Row}
\DeclareMathOperator{\nullspace}{Nul}
\DeclareMathOperator{\rank}{rank}
\DeclareMathOperator{\proj}{proj}
% -------------------------------------------------------

% ------- Set up circle -------
\newcommand{\bubble}{\tikz[baseline=-0.625ex]\draw[radius=0.5em, fill=white] (0,0) circle;}
\newcommand{\dbubble}{\tikz[baseline=-0.625ex]\draw[radius=0.5em, fill=black] (0,0) circle;}
% -------------------------------------------------------

\makeatletter
\setlength{\@fptop}{0pt}
\makeatother

\begin{document}
\maketitle

\section{\centering T/F Questions}
\begin{newquestion}
	For each statement, mark the statement \emph{true} if it's \emph{always} true, or \emph{false} otherwise.
	\begin{table}[!h]
		\centering
		\begin{tabular}{c@{\hspace{2.5em}}cp{0.75\textwidth}}
			true & false & \\
			\specialrule{\lightrulewidth}{0.25em}{1.25em}
			\bubble & \bubble & If \(A\) is \(4 \times 3\) with \(\dim(\rowspace (A)^\perp) = 2\), then \(\rank A = 2\). \\[1em]
			\bubble & \bubble & If \(\vec{x}\) belongs to \((\colspace A)^\perp\), then \(A^T\vec{x} = \vec{0}\).\\[1em]
			\bubble & \bubble & Given a matrix \(A\) is orthogonal, then it is also invertible.\\[1em]
			\bubble & \bubble & If a matrix \(A\) has a \(QR\) factorization, then \(\colspace A = \colspace Q\).\\[1em]
			\bubble & \bubble & Given a subspace \(W\) of \(\mathbb{R}^3\), any vector \(\vec{v} \in \mathbb{R}^3\) can be expressed as a sum of a
						      vector in \(W\) and \(W^\perp\). \\[2em]
			\bubble & \bubble & The Gram-Schmidt process applied on a basis of a subspace \(W\) does not necessarily yield an orthonormal basis for \(W\).\\[1em]
			\bubble & \bubble & If \(\nullspace(A^T) = \{\vec{0}\}\), then \(\vec{x} \mapsto A\vec{x}\) defines an onto linear transformation. \\[1em]
			\bubble & \bubble & If \(A\) is diagonalizable and similar to \(B\), then \(B\) is also diagonalizable. \\[1em]
			\bubble & \bubble & \(\proj_W \vec{y} = \vec{0}\) if and only if \(\vec{y} \in W^\perp\).\\[1em]
			\bubble & \bubble & If \(A \in \mathbb{R}^{2 \times 2}\) and has a complex eigenvalue \(\lambda_1\), then 
						      \(\lvert \lambda_1 \rvert = \lvert \lambda_2 \rvert\), where \(\lambda_2\) is the second eigenvalue of \(A\).\\[2em]
			\bubble & \bubble & For any orthogonal matrix \(A\), \(\lvert \det(A) \rvert = 1\).\\[1em]
			\bubble & \bubble & If \(A\) is diagonalizable with the diagonalization \(A = PDP^{-1}\), then we have \(\det(A) = \det(D)\). \\[1.75em]
			\bottomrule
		\end{tabular}
	\end{table}
\end{newquestion}

\section{\centering Possible/Impossible Questions}
\begin{newquestion}
	Mark each statement as \textit{possible} if it could ever be true/possible, or \textit{impossible} otherwise.
	\begin{table}[!h]
		\centering
		\begin{tabular}{c@{\hspace{2.5em}}cp{0.75\textwidth}}
			possible & impossible & \\
			\specialrule{\lightrulewidth}{0.25em}{1.25em}
			\bubble & \bubble & \(A^TA = AA^T = I\) and \(\lambda = 0\) is an eigenvalue of \(A\). \\[0.75em]
			\bubble & \bubble & An equation \(A\vec{x} = \vec{b}\) that doesn't have a least-squares solution. \\[0.75em]
			\bubble & \bubble & An onto linear transformation \(T \colon \mathbb{R}^3 \to \mathbb{R}^2\) whose standard matrix is \(A\) and 
						      \(\dim(\nullspace(A^T)^\perp)\) = 1.\\[1.75em]
			\bubble & \bubble & A \(3 \times 3\) real matrix with only complex eigenvalues. \\[0.75em]
			\bubble & \bubble & A \(3 \times 2\) matrix whose column space is \(W\) and \(W^\perp\) is a line. \\[0.75em]
			\bubble & \bubble & \(U \in \mathbb{R}^{n \times p}\)  with orthonormal columns and \(p < n\) with \(UU^T\vec{x} = \vec{x}\) for all \(\vec{x} \in \mathbb{R}^n\). \\[0.75em]
			\bubble & \bubble & Two vectors \(\vec{v} \ne \vec{u}\) but \(\proj_W \vec{v} = \proj_W \vec{u}\), where \(W\) is a subspace of 
						     \(\mathbb{R}^n\).\\[0.75em]
			\bubble & \bubble & Eigenvalue \(\lambda\) of \(A\) with algebraic multiplicity 1 and geometric multiplicity 2.\\[0.75em]
			\bottomrule
		\end{tabular}
	\end{table}
\end{newquestion}

\section{\centering Short/Long Answer Questions}
\begin{newquestion}
	For each question, provide your work and answer.
	\begin{enumerate}
		\item \(\mathcal{W} = \vspan \left\{\!\begin{bmatrix}1 \\ 0 \\ 2\end{bmatrix}, \begin{bmatrix}-1 \\ 1 \\ 3\end{bmatrix}\!\right\}\). What is a basis
			for \(\mathcal{W}^\perp\)?
		\item Express \(\vec{x} = \begin{bmatrix}2 \\ 6\end{bmatrix}\) as a sum of a vector in \(\mathcal{U} = \vspan\{\vec{u}\}\) and a vector in 
			\(\mathcal{U}^\perp\), where \(\vec{u} = \begin{bmatrix}7 \\ 1\end{bmatrix}\).
		\item Given vectors \(\vec{u} = \begin{bmatrix}1 \\ 5 \\ -1\end{bmatrix}\), \(\vec{v} = \begin{bmatrix}-1 \\ 1 \\ 4\end{bmatrix}\), and
			\(\vec{w} = \begin{bmatrix}2 \\ 2 \\ 0\end{bmatrix}\), answer the following questions.
				\begin{enumerate}
					\item Find the distance between \(\vec{u}\) and \(\vec{v}\).
					\item Do vectors \(\{\vec{v}, \vec{u}, \vec{w}\}\) form an orthogonal set?
					\item What would be an orthogonal basis for \(\mathcal{S} = \vspan\{\vec{u}, \vec{w}\}\)?
				\end{enumerate}
		\item Find a 2-by-3 matrix \(A\) such that \((\colspace A)^\perp = \vspan\left\{\!\begin{bmatrix}2 \\ 2\end{bmatrix}\!\right\}\).
			\item If \(\vec{u}_1 = \begin{bmatrix}1 \\ 0 \\ 1\end{bmatrix}, \vec{u}_2 = \begin{bmatrix}1 \\ 1 \\ -1\end{bmatrix}, 
			      \vec{u}_3 = \begin{bmatrix}-1 \\ 2 \\ 1\end{bmatrix}\), and \(\mathcal{W} = \vspan \{\vec{u}_1, \vec{u}_2\}\):
			\begin{enumerate}
				\item What is the closest vector in \(\mathcal{W}\) to \(\begin{bmatrix}1 \\ 5 \\ 1\end{bmatrix}\)?
				\item Find the coordinate vector of \(\begin{bmatrix}3 \\ 1 \\ 1\end{bmatrix}\) under the basis formed by 
				\(\mathcal{B} = \{\vec{u}_1, \vec{u}_2\}\).
				\item Find a vector \(\vec{x} \in \mathbb{R}^3\) such that:
					\begin{enumerate}
						\item \(\vec{x}\) is orthogonal to \(\vec{u}_3\).
						\item \(\proj_{\mathcal{L}_1} \vec{x} = 5\vec{u}_1,\) where \(\mathcal{L}_1 = \vspan \{\vec{u}_1\}\)
						\item \(\proj_{\mathcal{L}_2} \vec{x} = 3\vec{u}_2,\) where \(\mathcal{L}_2 = \vspan\{\vec{u}_2\}\)
					\end{enumerate} 
			\end{enumerate}
		\item If \(\|\vec{v}\| = \|\vec{u}\| = 1\) and \(\vec{v} \cdot \vec{u} = 0\), then what is \(\|2\vec{v} + 3\vec{u}\|\)?
		\item If \(A \in \mathbb{R}^{16 \times 15}\) and \(\dim(\nullspace(A^T)) = 12\), what is \(\dim(\nullspace(A)^\perp) + \dim(\colspace(A))\)?
		\item Use the \(QR\) decomposition of the matrix \(A = QR = \displaystyle \frac{1}{\sqrt{2}}\begin{bmatrix}1 & 1 \\ -1 & 1\end{bmatrix}\begin{bmatrix}2 & 2 \\ 0 & 1\end{bmatrix}\) to find the least square solution to \(A\vec{x} = \begin{bmatrix}2\sqrt{2} \\ \sqrt{2}\end{bmatrix}\).
		\item Find the value(s) for the parameters \(\alpha, \beta\) that'll best fit the curve \[y = \alpha x^2 + \beta\] to the data points (-1, 1), (1, 3), (1, 2),
			(1, 1).
		\item Find the eigenvalues of the matrix below. Then, find an invertible matrix \(P\) and a diagonal matrix \(D\) such that \(A = PDP^{-1}\). 
			\[A = \begin{bmatrix}1 & 0 & 0 \\ 1 & 2 & 0 \\ -1 & 0 & 2\end{bmatrix}.\]
	\end{enumerate}
\end{newquestion}

\section{Solutions}

\subsection{Solutions to T/F Questions}
	\begin{table}[!h]
		\centering
		\begin{tabular}{c@{\hspace{2.5em}}cp{0.75\textwidth}}
			true & false & \\
			\specialrule{\lightrulewidth}{0.25em}{1.25em}
			\bubble & \dbubble & If \(A\) is \(4 \times 3\) with \(\dim(\rowspace (A)^\perp) = 2\), then \(\rank A = 2\). \\[1em]
			\dbubble & \bubble & If \(\vec{x}\) belongs to \((\colspace A)^\perp\), then \(A^T\vec{x} = \vec{0}\).\\[1em]
			\dbubble & \bubble & Given a matrix \(A\) is orthogonal, then it is also invertible.\\[1em]
			\dbubble & \bubble & If a matrix \(A\) has a \(QR\) factorization, then \(\colspace A = \colspace Q\).\\[1em]
			\dbubble & \bubble & Given a subspace \(W\) of \(\mathbb{R}^3\), any vector \(\vec{v} \in \mathbb{R}^3\) can be expressed as a sum of a
						      vector in \(W\) and \(W^\perp\). \\[2em]
			\dbubble & \bubble & The Gram-Schmidt process applied on a basis of a subspace \(W\) does not necessarily yield an orthonormal basis for \(W\).\\[2em]
			\dbubble & \bubble & If \(\nullspace(A^T) = \{\vec{0}\}\), then \(\vec{x} \mapsto A\vec{x}\) defines an onto linear transformation. \\[1em]
			\dbubble & \bubble & If \(A\) is diagonalizable and similar to \(B\), then \(B\) is also diagonalizable. \\[1em]
			\dbubble & \bubble & \(\proj_W \vec{y} = \vec{0}\) if and only if \(\vec{y} \in W^\perp\).\\[1em]
			\dbubble & \bubble & If \(A \in \mathbb{R}^{2 \times 2}\) and has a complex eigenvalue \(\lambda_1\), then 
						      \(\lvert \lambda_1 \rvert = \lvert \lambda_2 \rvert\), where \(\lambda_2\) is the second eigenvalue of \(A\).\\[2em]
			\dbubble & \bubble & For any orthogonal matrix \(A\), \(\lvert \det(A) \rvert = 1\).\\[1em]
			\dbubble & \bubble & If \(A\) is diagonalizable with the diagonalization \(A = PDP^{-1}\), then we have \(\det(A) = \det(D)\). \\[1.75em]
			\bottomrule
		\end{tabular}
	\end{table}

\subsection{Solutions to Possible/Impossible Questions}
	\begin{table}[!h]
		\centering
		\begin{tabular}{c@{\hspace{2.5em}}cp{0.75\textwidth}}
			possible & impossible & \\
			\specialrule{\lightrulewidth}{0.25em}{1.25em}
			\bubble & \dbubble & \(A^TA = AA^T = I\) and \(\lambda = 0\) is an eigenvalue of \(A\). \\[0.75em]
			\bubble & \dbubble & An equation \(A\vec{x} = \vec{b}\) that doesn't have a least-squares solution. \\[0.75em]
			\bubble & \dbubble & An onto linear transformation \(T \colon \mathbb{R}^3 \to \mathbb{R}^2\) whose standard matrix is \(A\) and 
						      \(\dim(\nullspace(A^T)^\perp)\) = 1.\\[1.75em]
			\bubble & \dbubble & A \(3 \times 3\) real matrix with only complex eigenvalues. \\[0.75em]
			\dbubble & \bubble & A \(3 \times 2\) matrix whose column space is \(W\) and \(W^\perp\) is a line. \\[0.75em]
			\bubble & \dbubble & \(U \in \mathbb{R}^{n \times p}\)  with orthonormal columns and \(p < n\) with \(UU^T\vec{x} = \vec{x}\) for all \(\vec{x} \in \mathbb{R}^n\). \\[0.75em]
			\dbubble & \bubble & Two vectors \(\vec{v} \ne \vec{u}\) but \(\proj_W \vec{v} = \proj_W \vec{u}\), where \(W\) is a subspace of 
						     \(\mathbb{R}^n\).\\[0.75em]
			\bubble & \dbubble & Eigenvalue \(\lambda\) of \(A\) with algebraic multiplicity 1 and geometric multiplicity 2.\\[0.75em]
			\bottomrule
		\end{tabular}
	\end{table}


\subsection{Solutions to Short/Long Answer Questions}
\begin{enumerate}
	\item[\textbf{Q1}] \(\left\{\begin{bmatrix}-2 \\ -5 \\ 1\end{bmatrix}\right\}\). Any scalar multiple of this vector works.
	
	\item[\textbf{Q2}] We do this by expressing \(\vec{x}\) as \[\vec{x} = \underbrace{\proj_\mathcal{U}\vec{x}}_{\in \mathcal{U}} + \underbrace{\vec{z}}_{\in \mathcal{U}^\perp}.\] In this case, 
				   \(\proj_\mathcal{U}\vec{x} = \displaystyle \frac{\vec{u} \cdot \vec{x}}{\vec{u} \cdot \vec{u}}\vec{u} = \frac{2}{5}\vec{u}\). So
				   \(\vec{z} = \vec{x} - \proj_\mathcal{U} \vec{x} = 4/5\begin{bmatrix}-1 \\ 7\end{bmatrix}.\)
	
	\item[\textbf{Q3}]
			\begin{enumerate}
				\item This is \(\|\vec{v} - \vec{u}\| = \|\vec{u} - \vec{v}\| = \sqrt{45}\).
				\item No, as \(\vec{u} \cdot \vec{w} \ne 0\).
				\item We apply Gram-Schmidt here. Let \(\vec{v}_1 = \vec{u}\). Now, \[\vec{v}_2 = \vec{w}  - \proj_{\vec{u}} \vec{w} 
					= 1/9\begin{bmatrix}14 \\ -2 \\ 4 \end{bmatrix}.\] So \(\{\vec{v}_1, \vec{v}_2\}\) is an orthogonal basis for this subspace.
			\end{enumerate}
			
	
	\item[\textbf{Q4}] Multiple answers. One such answer is \[\begin{bmatrix}1 & 1 & 1 \\ -1 & -1 & -1\end{bmatrix}.\] The most general answer is in the form 
			\[\begin{bmatrix}a & ca & da \\ -a & -ca & -da\end{bmatrix},\] where \(c, d \in \mathbb{R}.\) Note that the question is simply asking for
			a matrix whose columns are orthogonal to \(\begin{bmatrix}2 \\ 2\end{bmatrix}\). We know that the columnspace is of dimension 1 because
			of the fact that \(\dim((\colspace A)^\perp) = 1\). This means, you only need to find one vector orthogonal to 
			\(\begin{bmatrix}2 \\ 2\end{bmatrix}\) and the rest has to be a scalar multiple of that.

	\item[\textbf{Q5}]  Answers are below.
			\begin{enumerate}
				\item The closest vector to in \(\mathcal{W}\) to \(\vec{v} = \begin{bmatrix}1 \\ 5 \\ 1\end{bmatrix}\) is \(\proj_\mathcal{W} \vec{v}.\)
				In this case, that's 
				\[\proj_\mathcal{W} \vec{v}
				  = \frac{\vec{v} \cdot \vec{u}_1}{\vec{u}_1 \cdot \vec{u}_1}\vec{u}_1 
				   + \frac{\vec{v} \cdot \vec{u}_2}{\vec{u}_2 \cdot \vec{u}_2}\vec{u}_2 = \begin{bmatrix}8/3 \\ 5/3 \\ -2/3\end{bmatrix}\]
				\item Since the vector is already in \(\mathcal{W}\), finding it's projection amounts to finding a linear combination of \(\{\vec{u}_1, \vec{u}_2\}\)
					that'll yield \(\vec{v} = \begin{bmatrix}3 \\ 1 \\ 1\end{bmatrix}\). That is 
					\[\proj_\mathcal{W} \vec{v} = \frac{\vec{u}_1 \cdot \vec{v}}{\vec{u}_1 \cdot \vec{u}_1}\vec{u}_1 + \frac{\vec{u}_2 \cdot \vec{v}}{\vec{u}_2 \cdot \vec{u}_2}\vec{u}_2 = \begin{bmatrix}\vec{u}_1 & \vec{u}_2\end{bmatrix}\begin{bmatrix}\frac{\vec{u}_1 \cdot \vec{v}}{\vec{u}_1 \cdot \vec{u}_1}\\[0.5em] \frac{\vec{u}_2 \cdot \vec{v}}{\vec{u}_2 \cdot \vec{u}_2}\end{bmatrix}.\] The final vector is the coordinate of \(\vec{v}\) under the new basis. 
					Computing the values will show that the coordinate is \(\begin{bmatrix}2 \\ 1\end{bmatrix}\).
				\item Note that \(\{\vec{u}_1, \vec{u}_2, \vec{u}_3\}\) forms an orthogonal set. Therefore, if \(\mathcal{W} =  \vspan\{\vec{u}_1, \vec{u}_2\}\),
					 then \(\mathcal{W}^\perp = \vspan\{\vec{u}_3\}\).\footnote{This is shown in the textbook, and it's easy to verify for yourself.} Since we want
					 the vector to be orthogonal to \(\vec{u}_3\), it should live in \(\mathcal{W}\). Since the vector already lives in \(\mathcal{W}\), then the weights of the 
					projection are just its new coordinates under the basis formed by \(\mathcal{B} = \{\vec{u}_1, \vec{u}_2\}.\) (If you tell me the
					projection, you're also telling me the coordinates in this case!) This alone is enough to tell us that \(\vec{x} = 5\vec{u}_1 + 3\vec{u}_2.\)
				\end{enumerate}

	\item[\textbf{Q6}] \(\|2\vec{v} + 3\vec{u}\|^2 = 4\|\vec{v}\|^2 + 9\|\vec{u}\|^2 = 13.\) We've used the Pythagorean Theorem. 
			So \(\|2\vec{v} + 3\vec{u}\| = \sqrt{13}\).
	\item[\textbf{Q7}] \((\nullspace A)^\perp = \rowspace A,\) and \(\dim(\rowspace A) = \dim(\colspace A) = 4.\) So the sum is \textbf{8}. 

	\item[\textbf{Q8}] \[R\vec{\hat{x}} = Q^T\vec{b} \implies \begin{bmatrix}2 & 2 \\ 0 & 1\end{bmatrix}\vec{\hat{x}}
								 = \sqrt{2} \cdot \frac{1}{\sqrt{2}} \begin{bmatrix}1 & - 1 \\ 1 & 1\end{bmatrix}
								 \begin{bmatrix}2 \\ 1\end{bmatrix} = \begin{bmatrix}1 \\ 3\end{bmatrix}.\]

					Row-reduce the augmented matrix to reach to the solution. 
					\[\left(\!\!\begin{array}{cc|c}2 & 2 & 1 \\ 0 & 1 & 3\end{array}\!\!\right) 
					  \sim \left(\!\!\begin{array}{cc|c}2 & 0 & -5 \\ 0 & 1 & 3\end{array}\!\!\right)
					  \sim \left(\!\!\begin{array}{cc|c}1 & 0 & -5/2 \\ 0 & 1 & 3\end{array}\!\!\right).\]
					We get that \(\vec{\hat{x}} = \begin{bmatrix}-5/2 \\ 3\end{bmatrix}\).

	\item[\textbf{Q9}] Follow the typical procedure for solving these types of problems.
					\[\begin{aligned}
						1 &= \alpha (-1)^2 + \beta &&\implies 1 = \alpha + \beta \\
						3 &= \alpha (1)^2  + \beta &&\implies 3 = \alpha + \beta \\
						2 &= \alpha (1)^2  + \beta &&\implies 2 = \alpha + \beta \\
						1 &= \alpha (1)^2  + \beta &&\implies 1 = \alpha + \beta \\
					\end{aligned}\]
					Something weird is happening here. The matrix equation associated with this system looks like 
					\[\begin{bmatrix}1 & 1 \\ 1 & 1 \\ 1 & 1 \\ 1 & 1\end{bmatrix}\begin{bmatrix}\alpha \\ \beta\end{bmatrix} 
					  = \begin{bmatrix}1 \\ 3 \\ 2 \\ 1\end{bmatrix}\]
					\textbf{The columns do not form a linearly independent set, so we'll have an infinite number of least-square solutions for 
					\(\alpha, \beta\) where they all yield the same amount of error}. Moving on,
					\[A^TA = \begin{bmatrix}1 & 1 & 1 & 1 \\ 1 & 1 & 1 & 1\end{bmatrix}
					\begin{bmatrix}1 & 1 \\ 1 & 1 \\ 1 & 1 \\ 1 & 1\end{bmatrix} = \begin{bmatrix}4 & 4 \\ 4 & 4\end{bmatrix};\]

					\[A^T\vec{b} = \begin{bmatrix}1 & 1 & 1 & 1 \\ 1 & 1 & 1 & 1\end{bmatrix}\begin{bmatrix}1 \\ 3 \\ 2 \\ 1\end{bmatrix}
					 = \begin{bmatrix}7 \\ 7\end{bmatrix}.\]

					Now, we solve the normal equations \(A^TA\vec{\hat{x}} = A^T\vec{b}.\) 
					\[\left[\!\!\begin{array}{cc|c}4 & 4 & 7 \\ 4 & 4 & 7\end{array}\!\!\right]
					  \sim \left[\!\!\begin{array}{cc|c}4 & 4 & 7 \\ 0 & 0 & 0\end{array}\!\!\right] 
					  \sim \left[\!\!\begin{array}{cc|c}1 & 1 & 7/4 \\ 0 & 0 & 0\end{array}\!\!\right].\]
					We find that \[\begin{bmatrix}\alpha \\ \beta\end{bmatrix} = \begin{bmatrix}7/4 \\ 0\end{bmatrix} 
					+ \beta\begin{bmatrix}-1 \\ 1\end{bmatrix},\] where \(\beta\) is the free variable. Each choice of \(\alpha, \beta\) obtained from the
					parametric vector equation yields a curve with the same amount of error.

	\item[\textbf{Q10}] There are multiple answers. We present ours here.
					\[A = \begin{bmatrix}1 & 0 & 0 \\ -1 & 1 & 0 \\ 1 & 0 & 1\end{bmatrix}\begin{bmatrix}1 & 0 & 0 \\ 0 & 2 & 0 \\ 0 & 0 & 2\end{bmatrix}\begin{bmatrix}1 & 0 & 0 \\ -1 & 1 & 0 \\ 1 & 0 & 1\end{bmatrix}^{-1}.\]
\end{enumerate}

\end{document}